\subsection{Optimizers for large-batch training}

Prior work has explored optimizers designed specifically for
large-batch training, including LARS~\citep{you2017large} and
LAMB~\citep{you2019large}.  It is instructive to consider these prior
findings in light of the scaling laws from our paper.
%
In particular, both original BERT~\citep{devlin2019bert} and the LAMB
replication were trained on 85.2B tokens.  Applying our fitted
$\bcrit$ power law over $D$=85.2B, we obtain an estimated $\bcrit$ of
about 12M tokens.  Original BERT was trained for 90\% of steps with a
batch size of 65K tokens (512 sequences of length 128).  LAMB
increased the batch size to 4M tokens (32K sequences), justifying
their claim, ``BERT training can be reduced from 3 days to just 76
minutes'' \citep{you2019large}.  However, based on the predicted
$\bcrit$ of 12M, batch size 4M is still well within the expected range
of efficient batch sizes. Moreover, the LAMB paper later notes, ``we
did not observe any speedup by increasing the batch size from 65536 to
131072 [sequences, or 16.8M tokens].''  In other words, they reach the
point of diminishing return \emph{exactly where $B$ exceeds our
predicted $\bcrit$}.

It is likely that some optimization issues solved by LAMB (to enable
stable large-batch training) are solved other ways in modern LLM
training setups, via, e.g., gradient clipping, pre-LayerNorm
placement, better initialization and stability control through $\mup$,
etc.  Scaling $\lambda$ rather than $\eta$ with $B$, as we propose,
further supports stable, efficient training. However, gradient
redundancy imposes an inherent limit on useful batch sizes, ensuring
critical batch size remains relevant.

\subsection{Critical batch size}\label{sec:app_bcrit}

Observations of critical batch size have previously been related to
data complexity~\citep{golmant2018computational}, loss
curvature~\citep{ma2018power,zhang2019algorithmic}, and model
architecture~\citep{shallue2019measuring}.

\citet{merrill2025critical} define $\bcrit$ as the largest $B$ such
that loss does not degrade by more than a fixed fraction $\epsilon$
from the $\bopt$ setting.  They measure this $\bcrit$ instantaneously
throughout training, by repeatedly branching from a checkpoint with
different $B$ settings and assessing the impact on loss.

Recent work also defines $\cbs$ in terms of how $\eta$ scales with
$B$~\citep{filatov2024time,li2024surge}; unlike our work, these recent
studies use a constant learning rate schedule and no weight decay.

We follow \citet{mccandlish2018empirical}'s definition of $\bcrit$
(\cref{def:cbs}).
%
Given various theoretical assumptions, \citet{mccandlish2018empirical}
derived a direct equivalence between $\bcrit$ and what they call the
\emph{gradient noise scale} (GNS): the variation of the gradients
between different training examples.
%
However, they noted that the GNS ``accurately predicts the largest
usable batch size (at the order of magnitude level),'' which is below
the level of precision needed for large-scale training.
%
\citet{merrill2025critical} recently found ``the gradient noise scale
underestimates the CBS [i.e., $\bcrit$].''
%
This lack of precision may be why, in Kaplan et al's original scaling
laws paper~\citep{kaplan2020scaling}, they note that, ``although the
critical batch size roughly matches the gradient noise scale, we are
using a direct [empirical] measurement of $\bcrit$.'' Our approach to
measuring $\bcrit$ (\cref{sec:estimating_bcrit}) similarly provides a
direct empirical measurement, but one that can be efficiently computed
with any learning rate schedule or optimizer.

\subsection{Detailed comparison with \citet{zhang2024how}}\label{sec:appendix_zhang}

Here we provide further comparison with the concurrent work
by \citet{zhang2024how}.
%
The primary point of distinction of our paper is that we conducted a
large-scale empirical study into the scaling of AdamW's weight decay
hyperparameter (including its scaling with $B$), ultimately deriving a
precise power law for the optimal AdamW timescale in
tokens-per-parameter.  \citet{zhang2024how} did not use weight decay.
Further, we also explored scaling of $\bopt$ in addition to $\bcrit$.
%
Beyond use of weight decay, further methodological differences in our
main experiments include that we used a longer context length (2048
vs.\ 512), a cleaner dataset (SlimPajama vs. C4), the $\mup$
parameterization, a decaying LR schedule (more on this below), and
that we tuned HPs at most $N$, $D$, $B$ (\cref{tab:train_steps}),
while \citet{zhang2024how} performed a HP sweep for a 151M model, and
re-used optimal values at other scales.
%
We now focus on differences in estimating and measuring the scaling of
$\bcrit$.

\paragraph{Estimating $\bcrit$ for a specific target loss}

Both our work and \citet{zhang2024how} require measuring, for
different batch sizes, how many training steps it takes to reach a
particular target loss.
%
Since the number of steps to reach that loss is not known \emph{a
priori}, it is inherently difficult to study $\bcrit$ when using a LR
decay schedule, where you must specify the number of steps in advance.
Using a constant LR (as was done in early work on
$\bcrit$~\citep{mccandlish2018empirical}) simply does not result in
competitive models~\citep{bergsma2025straight}.
%
Unfortunately, it is not feasible to \emph{search} for the precise
step count needed, i.e., by conducting full training runs with
different schedules/step budgets.

\citet{zhang2024how} creatively solve this issue by conducting a
single training run at a constant LR, while using weight averaging to
generate higher-quality checkpoints for evaluation.  With this
approach, they still ``need to frequently evaluate the model on a
holdout evaluation set'' \citep{zhang2024how}.

Given LR decay, as opposed to weight averaging, remains the standard
practice for current state-of-the-art LLMs, we independently developed
a different approach.  This led to the novel method described in our
paper.  In contrast with~\citet{zhang2024how}, we do not need to
frequently evaluate the model, as we instead fit a $B$-specific loss
power law through a few validation loss values
(\cref{sec:estimating_bcrit}).  Indeed, our approach may improve the
efficiency of \citet{zhang2024how}'s method, as it would obviate the
cost of continuous validation, which concerned them (see their section
``Evaluation data size and frequency'').

\paragraph{Estimating the $\bcrit$ power law}

Collecting $\bcrit$ data across multiple model scales and loss targets
is expensive.  \citet{zhang2024how} establish $\bcrit$ scaling in $D$
through three targeted experiments:
\begin{itemize}
\item measuring $\bcrit$ while scaling $N$ but leaving $D$ fixed to 3.07B
\item measuring $\bcrit$ while scaling $D$ but leaving $N$ fixed to 302M
\item measuring $\bcrit$ while scaling both $N$ and $D$ proportionally (at 20 TPP)
\end{itemize}
Interestingly, $\bcrit$ was found to only scale weakly in $N$, but
scale similarly whenever $D$ is scaled.  They then fit a power law to
their data points for the 302M-parameter model, obtaining the fit
\mbox{$\bcrit = 22.91 D^{0.47}$} (in tokens).

In comparison, we took a more brute-force approach, computing $\bcrit$
across many different $N$ and $D$ values, and ultimately fitting our
$\bcrit$ power law across multiple different model sizes and TPP
settings (\cref{fig:hook}, $\rightfig$).  Also,
unlike~\citet{zhang2024how}, we assessed the quality of fit via
computation of both $\Rtwo$ values and parameter quantiles via
bootstrapping (\cref{sec:results_cbs}).

Recall also that \citet{zhang2024how} used a different definition of
critical batch size.  Let us denote their quantity $\cbszhang$.  They
set $\cbszhang$ to be the $B$ such that the data required to reach a
loss target is $1.2 \times \dmin$ (i.e., $1.2 \times$ the data
required with $\bopt$).
%Their fitted power law is
%\mbox{$\cbszhang = 22.91 D^{0.47}$}, where $\cbszhang$ is in tokens.

We can use \cref{eq:s_extra} to align their fitted law with our own.
By this equation, we have:
\begin{align*}
  D &= \dmin(1 + \frac{\cbszhang}{\cbs}) \\
    &:= \dmin(1.2) \\
  \Rightarrow \frac{\cbszhang}{\cbs} &= 0.2 \\
  \Rightarrow \cbs &= 5 \cbszhang
\end{align*}
%
Thus, to convert their coefficient to our scale, we multiply it by 5,
and, dividing by the number of tokens in our sequences, obtain
$\cbszhang = 0.0559 D^{0.47}$.
%
The $\cbszhang$ coefficient (0.0559) is 19\% larger than our own
(0.0471), perhaps reflecting differences in training setup or data
quality (and worth investigating further in future work).
%
However, the exponents are quite similar (0.47 vs. 0.462), suggesting
that both works are independently measuring the same fundamental
scaling behavior.

We emphasize that $\bcrit$ directly reflects the fundamental limit to
data parallelism in training neural networks.  Given the significant
implications of $\bcrit$ scaling in $D$ rather than $C$ or $L$
(including those discussed in \cref{sec:balancing}), we note the
scientific and practical value in having different approaches
independently observe this same phenomenon.

\subsection{Hyperparameter scaling with $B$}

It has long been recognized that the optimal learning rate, $\etaopt$, scales
with $B$, with reports of both
linear~\citep{krizhevsky2014one,chen2016revisiting,smith2017bayesian,smith2018dont}
and square-root
scaling~\citep{hoffer2017train,you2019large,malladi2022sdes}.
%
Recent work has found $\etaopt$ to \emph{decrease} when $B
> \cbs$~\citep{li2024surge,filatov2024time}, which resonates with our
own findings (\cref{fig:tepochs}, $\rightfig$).
%
%\citet{li2024surge} look at learning rate scaling with Adam, they see an
%increase, then a decrease, then an increase again, but at many
%multiples of $\cbs$...
%- That Time Transfer paper also sees this with MUP
%
Since it is difficult to predict exactly how $\eta$ will scale with
$B$, studies of $\cbs$ have often done full HP sweeps at each
$B$~\citep{mccandlish2018empirical,shallue2019measuring}.
% \citet{zhang2024how} perform a HP sweep for a 151M model, and then
%re-use optimal values at other scales.

The only work we are aware of that specifically recommends scaling
weight decay with $B$ is \citet{loshchilov2017decoupled}, who suggest
$\lambda \propto \sqrt{B}$, though this rule is not evaluated
systematically. It is also important to note that
\citet{loshchilov2017decoupled} use the \emph{independent} form of
weight decay, where decay is applied independently of the learning
rate $\eta$, unlike common implementations such as AdamW in
PyTorch~\citep{wortsman2023small}. In these more typical
\emph{dependent} implementations, weight decay is scaled by $\eta$, so
any increase in $\eta$ with $B$ (e.g., $\eta \propto B$ or $\sqrt{B}$)
already increases the effective weight decay strength accordingly.

%399R: Note in Loshchilov & Hutter (2017) they use fully decoupled weight decay. Scaling it with the sqrt(B) may therefore simply be equivalent to scaling the learning rate by sqrt(B) in the more common not-fully-decoupled AdamW implementation (default in PyTorch).

\subsection{$\tepoch$ and effective learning rates}

The concept of effective learning rates, influenced by weight decay,
has been widely
discussed~\citep{van2017l2,hoffer2018norm,zhang2018three,chiley2019online,li2019exponential,wan2020spherical,kosson2023rotational,dangelo2024why}.
%
In its simplest form, the effective or \emph{intrinsic} LR is simply
$\eta\lambda$, but in these prior works, effective LRs typically
measure functional updates relative to weight magnitude, which is
particularly relevant for normalization-based networks.
%
Comparison of the effects of $\lambda$ vs.\ $\eta$ adjustments in the
context of LR decay schedules was explored
in~\citet{bergsma2025straight}.

The behavior of effective LRs (relative update sizes) over the course
of training has been studied comprehensively
by \citet{kosson2023rotational}, including comparing the effects of
increasing $\eta$ vs.\ increasing $\lambda$.
%
This work shows that higher $\eta$ values can cause large relative
updates early in training, which can destabilize training or require
longer warmups~\citep{kosson2024analyzing}.
%
High $\eta$ and low $\lambda$ can also lead to larger weight
norms~\citep{kosson2023rotational,dangelo2024why}, which also has a
destabilizing effect, particularly on low-precision training.
%
These effects may explain why we were able to achieve higher effective
LRs $\eta\lambda$ by tuning $\lambda$ rather than tuning $\eta$ with
$B$ (\cref{fig:tepochs}, $\middlefig$).

For a given dataset size $D$, the $\tepoch$ and the batch-normalized
effective LR $\frac{\eta\lambda}{B}$ are equivalent, and thus
effective LRs and the AdamW timescale can be viewed as different
perspectives on the AdamW optimization process.
